\section{Research Framework and Hypotheses}
\label{sec:research_framework}

The literature reviewed in this chapter supports a coherent framework linking institutional ownership to future downside risk. Institutional investors may hold overlapping or otherwise correlated positions, while common liquidity shocks can induce simultaneous trading and non-fundamental price pressure. A stock's vulnerability to this mechanism therefore depends not only on its aggregate institutional ownership, but also on the composition and structure of its investor base. Conventional ownership measures summarize selected aspects of that exposure, whereas a network representation preserves the relationships among managers and securities. Predictive models can then test whether this additional information improves out-of-sample forecasts and whether learning directly from the ownership graph provides an advantage over engineered network features.

This framework motivates five hypotheses organized around three dimensions: the information contained in ownership data, the ability of different models to exploit that information, and the economic relevance of the resulting predictions. The hypotheses are directional predictions to be evaluated empirically rather than maintained assumptions; failure to find an incremental benefit from ownership information or model complexity would therefore remain an informative result.

\begin{description}[leftmargin=2.2em, style=nextline]

      \item[H1 --- Conventional Ownership and Crowding.]
            Conventional institutional-ownership and crowding measures improve the out-of-sample prediction of stock-level downside risk relative to market characteristics alone.

      \item[H2 --- Ownership--Network Topology.]
            Engineered features describing ownership-network topology provide incremental predictive information beyond market characteristics and conventional ownership measures.

      \item[H3 --- Nonlinearity.]
            Nonlinear tabular models predict future downside risk more accurately than linear models estimated on the same feature sets.

      \item[H4 --- Graph Learning.]
            Graph-based models that learn directly from the manager--stock ownership network predict future downside risk more accurately than comparable tabular models, including those supplied with engineered network features.

      \item[H5 --- Economic Value.]
            The model-derived Fragility Score produces an economically meaningful out-of-sample ranking of subsequent downside risk and provides useful information for portfolio construction.

\end{description}

\noindent
Hypotheses H1 and H2 test the incremental information contained in conventional ownership measures and ownership-network topology. Hypotheses H3 and H4 examine whether greater model flexibility improves the extraction of information from a fixed feature set or directly from the ownership graph. Hypothesis H5 extends the evaluation beyond statistical accuracy to the economic significance of the resulting predictions. Chapter~\ref{chap:methodology} translates these hypotheses into explicit empirical comparisons, while Chapter~\ref{chap:results} reports the corresponding evidence.